\chapter{Trong Lớp, Mắt Em Ở Đâu?}
\markboth{Trong Lớp, Mắt Em Ở Đâu?}{Trong Lớp, Mắt Em Ở Đâu?}

\begin{center}
	\textit{\color{bookprimary}``Có mặt trong lớp không có nghĩa là đầu óc cũng có mặt ở đó.''}

	\textit{\color{bookprimary}``Being in the classroom does not mean your mind is there too.''}
\end{center}

\ngancach

\begin{tcolorbox}[colback=booksecondary!8,colframe=booksecondary!40,arc=5pt,title={\bfseries Góc Suy Nghĩ},fonttitle=\bfseries\color{booksecondary}]
Nhiều học sinh nghĩ điện thoại chỉ lấy mất thời gian giải trí.
\textit{(Many students think the phone only steals entertainment time.)}

Không.
\textit{(No.)}

Trong lớp, nó lấy thẳng khả năng nối những ý đang được giảng thành một mạch hiểu.
\textit{(In class, it directly steals the ability to connect the ideas being taught into one line of understanding.)}

Một lần cúi xuống tưởng chỉ mất vài giây.
\textit{(One glance downward seems to cost only a few seconds.)}

Nhưng cái bị đứt không chỉ là vài giây.
\textit{(But what breaks is not only a few seconds.)}

Nó là sợi chỉ đang nối em với bài học ngay lúc đó.
\textit{(It is the thread linking you to the lesson at that very moment.)}
\end{tcolorbox}

\section{Nghe Giảng Nửa Vời}
\begin{truyen}{Nghe Giảng Nửa Vời}{Half-Present Learning}
\chuhoa{C}{ó bạn để điện thoại dưới ngăn bàn, mắt vẫn nhìn lên như rất ngoan, nhưng ngón tay thì lướt nhẹ trong bóng tối.}
\textit{(A student keeps the phone under the desk, eyes still looking up as if very obedient, while the fingers quietly scroll in the dark.)}

Thầy gọi lên bảng.
\textit{(The teacher calls the student to the board.)}

Bạn ấy đứng đực ra vì không biết lớp đang ở bước nào.
\textit{(The student stands frozen, not knowing which step the class is on.)}

\teacherpair{Em tưởng em đang ăn gian được vài phút giải trí. Thật ra em đang tự bẻ gãy dòng suy nghĩ của chính mình. Học không chỉ là nghe từng câu rời. Học là giữ được một mạch tập trung đủ lâu để ý này nối sang ý khác. Điện thoại chen vào giữa mạch đó giống như có người cứ thò tay cắt đôi câu chuyện trong đầu em.}{``You think you are stealing a few minutes of entertainment. In fact, you are breaking your own chain of thought. Learning is not just hearing separate sentences. Learning means holding concentration long enough for one idea to connect to the next. When a phone cuts into that chain, it is like someone repeatedly slicing your thought in half.''}

Thầy nhìn cả lớp rồi nói thêm: ``Nhiều em không yếu trí thông minh. Các em yếu sức giữ chú ý. Mà người không giữ nổi chú ý thì rất khó dùng hết trí thông minh của mình.''
\textit{(The teacher looks at the whole class and adds: ``Many of you are not weak in intelligence. You are weak in holding attention. And a person who cannot hold attention will rarely be able to use their intelligence fully.'')}
\end{truyen}

\section{Chụp Bảng Thay Vì Hiểu Bài}
\begin{truyen}{Chụp Bảng Thay Vì Hiểu Bài}{Photographing the Board Instead of Understanding It}
\chuhoa{C}{ó học sinh không chép, không hỏi, chỉ giơ máy lên chụp bảng rồi yên tâm như thể kiến thức đã thuộc về mình.}
\textit{(A student neither writes nor asks, but simply takes a photo of the board and feels secure as if the knowledge now belongs to them.)}

Ảnh chụp giúp lưu thông tin.
\textit{(A photo helps store information.)}

Nhưng nó không tự biến thành hiểu biết.
\textit{(But it does not transform itself into understanding.)}

\teacherpair{Lưu được chưa chắc tiêu hóa được. Nhiều em chụp để khỏi phải vật lộn ngay lúc đó, rồi sau này cũng chẳng mở lại. Kiến thức bị đẩy khỏi đầu em sang thư viện ảnh, còn đầu em thì vẫn đói.}{``Storing something does not mean digesting it. Many students take a photo to avoid struggling in the moment, then never reopen it. The knowledge gets pushed from the mind into a photo library, while the mind itself stays hungry.''}
\end{truyen}

\section{Làm Bài Nhưng Chờ Thông Báo}
\begin{truyen}{Làm Bài Nhưng Chờ Thông Báo}{Doing the Work While Waiting for Notifications}
\chuhoa{C}{ó bạn đặt điện thoại úp xuống bàn lúc làm bài, không nhìn trực tiếp nhưng cả đầu óc vẫn chờ nó rung.}
\textit{(A student puts the phone face down while working, not looking directly at it yet mentally waiting for it to vibrate.)}

Thân xác có thể ngồi trước vở.
\textit{(The body may sit before the notebook.)}

Nhưng sự chú ý thì đang trực chiến ở nơi khác.
\textit{(But attention is stationed elsewhere.)}

\teacherpair{Không phải chỉ lúc em mở máy mới mất tập trung. Chỉ cần em để một phần đầu óc đứng gác chờ thông báo là sức học đã bị rút bớt rồi.}{``You do not lose focus only when you open the phone. The moment part of your mind stands guard waiting for a notification, your learning strength has already been drained.''}
\end{truyen}

\section{Hỏi Bạn Vì Mất Một Đoạn}
\begin{truyen}{Hỏi Bạn Vì Mất Một Đoạn}{Asking Friends Because You Missed One Piece}
\chuhoa{C}{ó bạn chỉ cúi xuống xem máy một lúc ngắn, ngẩng lên đã phải hỏi: ``Ủa nãy thầy đi tới đâu rồi?''}
\textit{(A student glances down at the phone for a short moment, looks up, and has to ask, ``Where did the teacher get to just now?'')}

Một đoạn mất đi kéo theo cả chuỗi sau đó khó hiểu hơn.
\textit{(One missed segment makes the rest of the sequence harder to understand.)}

\teacherpair{Trong học tập, có những thứ mất chỉ một chút mà hại rất lâu. Vì bài giảng có mạch. Em đứt ở một chỗ, nhiều khi những chỗ sau không còn đứng vững trong đầu em nữa.}{``In learning, some things are lost for a moment but hurt for a long time. A lesson has a flow. If you break at one point, later parts often fail to stand firmly in your mind.''}
\end{truyen}

\section{Mỗi Lần Rung Là Mỗi Lần Rơi}
\begin{truyen}{Mỗi Lần Rung Là Mỗi Lần Rơi}{Every Vibration Is Another Drop}
\chuhoa{C}{ó lớp để điện thoại trong cặp nhưng chưa tắt âm báo. Mỗi lần máy rung, mắt vài em lại chệch đi như bị giật dây.}
\textit{(A class keeps phones in bags without muting notifications. Each vibration pulls some eyes away as if by a string.)}

Đôi khi không cần mở máy, chỉ một tín hiệu cũng đủ làm đầu óc trượt khỏi chỗ đang đứng.
\textit{(Sometimes the phone need not even be opened; one signal alone is enough to make the mind slip.)}

\teacherpair{Đừng đánh giá thấp những cú rung nhỏ. Một tiết học không sụp vì một trận lớn. Nó bị khoét rỗng bằng nhiều mũi kim rất nhỏ như thế.}{``Do not underestimate small vibrations. A lesson rarely collapses in one great blow. It gets hollowed out by many tiny needles like that.''}
\end{truyen}

\section{Ghi Chép Đẹp Nhưng Đầu Rỗng}
\begin{truyen}{Ghi Chép Đẹp Nhưng Đầu Rỗng}{Beautiful Notes, Empty Understanding}
\chuhoa{C}{ó bạn ghi bài rất sạch, bút màu đủ hết, tiêu đề đẹp như in, nhưng khi hỏi lại thì không giải thích nổi ý chính.}
\textit{(A student keeps beautiful colorful notes, yet cannot explain the main idea when asked.)}

Đó là lúc việc học bị trượt từ hiểu sang trình bày.
\textit{(That is when learning slips from understanding into presentation.)}

\teacherpair{Hình thức không có tội. Nhưng nếu em chăm phần đẹp hơn phần hiểu, em đang biến việc học thành một màn dàn dựng dễ nhìn mà rỗng ruột.}{``Appearance is not guilty in itself. But if you care more for neatness than understanding, you are turning learning into a performance that looks good while hollow inside.''}
\end{truyen}

\section{Học Nhóm Mà Ai Cũng Phân Tán}
\begin{truyen}{Học Nhóm Mà Ai Cũng Phân Tán}{Study Groups Where Everyone Is Scattered}
\chuhoa{C}{ó nhóm ngồi học chung nhưng cứ vài phút lại có người kéo cả bàn lệch sang một video, một tin nhắn, một chuyện ngoài lề.}
\textit{(A study group sits together, but every few minutes someone drags the whole table toward a video, a message, or something unrelated.)}

Ngồi chung không tự động sinh ra tập trung.
\textit{(Sitting together does not automatically create focus.)}

Nếu cả nhóm đều yếu trước màn hình, sự phân tán sẽ lây rất nhanh.
\textit{(If the whole group is weak before screens, distraction spreads very quickly.)}

\teacherpair{Bạn học tốt không phải là người chỉ giỏi giảng. Còn là người biết giữ bầu khí học không bị rách bởi quá nhiều thứ lặt vặt.}{``A good study companion is not only someone who explains well. It is also someone who helps protect the learning atmosphere from being torn apart by trivial distractions.''}
\end{truyen}

\section{Cất Máy Xa Là Điểm Bắt Đầu}
\begin{truyen}{Cất Máy Xa Là Điểm Bắt Đầu}{Putting the Phone Far Away Is a Beginning}
\chuhoa{C}{ó học sinh thử để điện thoại ngoài tầm tay suốt một tiết và ngạc nhiên vì lần đầu nghe bài liền mạch hơn hẳn.}
\textit{(A student tries keeping the phone out of reach for one whole lesson and is surprised at how much more continuous the learning feels.)}

Nhiều khi cải thiện sự tập trung không bắt đầu từ động lực lớn.
\textit{(Sometimes better concentration does not begin with grand motivation.)}

Nó bắt đầu từ một thay đổi rất vật lý: kéo nguồn phá mạch ra xa mình.
\textit{(It begins with a physical change: moving the source of interruption farther away.)}

\teacherpair{Đừng coi thường những biện pháp đơn giản. Có khi trưởng thành bắt đầu bằng việc em chấp nhận bất tiện một chút để đầu óc mình được nguyên vẹn hơn.}{``Do not underestimate simple measures. Sometimes maturity begins when you accept a little inconvenience so your mind can remain more whole.''}
\end{truyen}

\begin{baihoc}
Điện thoại trong lớp không chỉ làm em mất tập trung.
\textit{(A phone in class does not just distract you.)}

Nó dạy đầu óc em một thói quen rất nguy hiểm.
\textit{(It teaches your mind a very dangerous habit.)}

Đó là cứ hễ có thứ gì hơi chậm hoặc đòi nghĩ là lập tức bẻ sang một kích thích khác.
\textit{(That habit is to jump to another stimulation the moment something becomes slow or demanding.)}
\end{baihoc}

\begin{tuhocnhanh}
1. Khi học, em bị đứt mạch bởi điện thoại bao nhiêu lần trong một tiết? \textit{(During one lesson, how many times does the phone break your chain of thought?)}

2. Em có đang lầm giữa việc ngồi trong lớp với việc thật sự học không? \textit{(Are you confusing sitting in class with actually learning?)}

3. Nếu cất điện thoại ra xa một tiết, em sợ điều gì mất nhất? \textit{(If you place the phone far away for one lesson, what do you fear losing most?)}
\end{tuhocnhanh}

\begin{tongketchuong}
Nhiều học sinh không học kém vì đầu óc quá yếu.
\textit{(Many students do not learn poorly because their minds are too weak.)}

Các em học kém vì sự chú ý của mình bị bán lẻ từng mẩu cho quá nhiều thứ rẻ tiền.
\textit{(They learn poorly because their attention is retail-sold in tiny pieces to too many cheap distractions.)}
\end{tongketchuong}

\ngancach